\subsection{Non-uniqueness of Spectral Reconstruction}
  \label{subsec:non-uniqueness}

  Reconstruction of an effective metric from spectral data is generically non-unique.
  This ambiguity is structural, not a consequence of approximation.

  Distinct relational configurations may share identical low-resolution
  spectral content within an admissible window and therefore yield the same
  effective metric.
  Conversely, a given configuration may admit multiple spectrally equivalent
  continuum embeddings.

  Because distance is defined operationally through spectral proximity,
  different filtering scales or embedding procedures can produce metrically
  equivalent but geometrically distinct descriptions.
  All such descriptions agree on observable properties within their common
  domain of validity.
  The effective metric thus represents an equivalence class compatible with
  the same spectral data.

  \paragraph{Flux conservation and radial scaling.}

    The emergence of a $1/r$ profile follows from conservation of relational
    relaxation flux in quasi-static regimes.
    The total flux through any closed relational surface surrounding a localized
    perturbation is conserved.

    In an effective three-dimensional projectable regime,
    flux conservation implies a $1/r^{2}$ decay of connectivity perturbations.
    Integration along admissible paths then yields an effective potential
    proportional to $1/r$.
    The Schwarzschild factor $1 - r_s/r$ thus arises from isotropic
    flux conservation rather than from an imposed force law.

  \paragraph{Origin of the $1/r$ profile.}

    Mass corresponds operationally to a localized inhibition of admissible
    relaxation in the relational substrate.
    This reduces the density of optimal relational paths.

    In the continuum regime, this reduction acts as a source term for the
    effective scalar Laplacian governing metric deformations.
    In three spatial dimensions, the corresponding Poisson equation has the
    unique fundamental solution proportional to $1/r$.
    The Schwarzschild scaling therefore follows from flux conservation in an
    isotropic effective geometry, not from an assumed gravitational law.
